\documentclass{exam}

\newif\ifprintselected   % required by exam_config's \ansbox

%%%%%%%%%%%%%%%%%%%%%%%%%%%%%%%%%%%%%%%%%%%%%%%%%%%%%%
%%% SET THESE IF VARIABLES
%%%%%%%%%%%%%%%%%%%%%%%%%%%%%%%%%%%%%%%%%%%%%%%%%%%%%%
% v3 of the midterm packet, for Fall 2027. v2's open-ended brief, plus what
% the Fall 2026 class needed: graded checkpoints, a reference antenna to check
% their setup against, a coax-stub match in place of the discrete L-network
% (which worked for nobody), and ESD protection on the dipole itself.
% Instructor-side notes, including the VNA port protector, are in
% project/MIDTERM_INSTRUCTOR_NOTES.md and deliberately not in this packet.
\ifdefined\thisTermIn
	\def\thisterm{\thisTermIn}
\else
	\def\thisterm{Fall 2027}
\fi
\def\thisclass{ECE 444}

\def\fOne{915}

% Due date and weights. Both are quoted from the syllabus -- change there first.
\def\duedate{Lesson 20, 2359}
% Checkpoint lessons. Each is due at the start of class.
\def\cpOne{11}
\def\cpTwo{14}
\def\cpThree{16}
\def\cpFour{18}
\def\projweight{30\%}
\def\midweight{40\%}

%%%%%%%%%%%%%%%%%%%%%%%%%%%%%%%%%%%%%%%%%%%%%%%%%%%%%%
%%% END SET THESE IF VARIABLES
%%%%%%%%%%%%%%%%%%%%%%%%%%%%%%%%%%%%%%%%%%%%%%%%%%%%%%
\input{myPackages_gr}
\usepackage[hidelinks, linkcolor=red]{hyperref}
\input{myShortcuts}
% ECE 444 document fonts.
%
% Body text is Barlow (SIL OFL; the faces and the license are in
% latex/fonts/Barlow/). Barlow is not in TeX Live, so it is loaded as an
% OpenType face through fontspec -- which means the document must be built
% with lualatex, not pdflatex. build_practice.sh and build_project.sh do; the
% lab packets build in ECE444-faculty from its vendored copy of this file.
%
% Compiling with pdflatex anyway is not an error: the guard below falls back
% to the LaTeX default so the source still compiles under any engine and can
% be syntax-checked. The resulting PDF just is not the one we ship.
%
% Paths are relative to latex/, which is where the build scripts run
% (each does cd "$(dirname "$0")"). Loading by path rather than by installed
% family name means no system font install and no fontconfig, so a fresh
% clone builds a byte-comparable PDF anywhere.
%
% MATH IS DELIBERATELY LEFT IN COMPUTER MODERN. Barlow has no math companion,
% and CM math under sans body text is already the house pairing: the decks set
% words in Source Sans Pro and let MathJax render symbols in its CM-derived
% face. Moving math to a sans as well would mean unicode-math plus a sans math
% font (Fira Math being the plausible one) and re-checking every equation in
% the course for symbol coverage. That is a separate decision, not a side
% effect of this one.

\usepackage{iftex}
\ifLuaTeX
	\usepackage{fontspec}
	\defaultfontfeatures{Ligatures=TeX}
	\setmainfont{Barlow}[
		Path           = fonts/Barlow/,
		Extension      = .ttf,
		UprightFont    = *-Regular,
		ItalicFont     = *-Italic,
		BoldFont       = *-Bold,
		BoldItalicFont = *-BoldItalic,
	]
	\setsansfont{Barlow}[
		Path           = fonts/Barlow/,
		Extension      = .ttf,
		UprightFont    = *-Regular,
		ItalicFont     = *-Italic,
		BoldFont       = *-SemiBold,
		BoldItalicFont = *-BoldItalic,
	]
	% microtype earns its place here specifically because of Barlow. Its
	% glyphs are narrower than Computer Modern's, so justified prose sets
	% tighter and TeX has to squeeze interword space to fit; two paragraphs
	% in the L02/L03 sets went overfull on the first Barlow build. Protrusion
	% and font expansion absorb that -- both are lualatex features, so this
	% belongs inside the guard with everything else.
	\usepackage{microtype}

	% Mono stays Latin Modern: Barlow ships no monospaced face, and verbatim
	% and code need one that actually is fixed-pitch.
\fi

\geometry{margin=0.75in}   % holds the packet to one page
\setlength{\parindent}{0pt}
\setlength{\parskip}{4pt}

\input{exam_config}
\noprintanswers
\printselectedfalse

% Compact run-in headings: \section's spacing alone pushed v2 onto a second page.
\newcommand{\head}[1]{\par\vspace{6pt}{\large\bfseries #1}\par\vspace{-2pt}}

\firstpageheader{}{}{}
\runningheader{\thisclass \hspace{1pt} -- \thisterm }{Midterm Project}{\thepage\hspace{1pt} of\hspace{1pt} \numpages}
\firstpagefooter{}{}{}
\runningfooter{}{}{}
\runningheadrule

\begin{document}
\pagestyle{headandfoot}

\begin{center}
USAF ACADEMY \\
DEPARTMENT OF ELECTRICAL AND COMPUTER ENGINEERING \\
MIDTERM PROJECT (assigned Lesson 9) - Dipole Design, Matching, and Measurement \\
\thisterm \\[2mm]
\end{center}

\textbf{Due \duedate.} Individual. Worth \projweight\ of the course grade and
\midweight\ of your midterm grade. There is no resubmission.

\head{The Assignment}

Build a dipole, make it work at the frequency it was built for, then make the
same antenna work at a frequency it was not built for. Measure what each
approach does to the impedance and to the pattern, and write it up.

\begin{enumerate}[nosep]
	\item Design and build a dipole for \fOne\ MHz.
	\item Tune it by trimming.
	\item Characterize it with S-parameters and input impedance.
	\item Measure its pattern in the chamber, co-polarized and cross-polarized.
	\item Pick a different frequency and measure the antenna's impedance there.
	      Leave the antenna as it is.
	\item Design and build a coax stub that matches it at that frequency.
	\item Measure the matched antenna's S-parameters and its pattern, co-polarized
	      and cross-polarized.
	\item Model the antenna you built in 4NEC at both frequencies and compare the
	      model against your measurements.
	\item Write a report detailing your design, your process, and your analysis.
\end{enumerate}

\textbf{Two build requirements.} The chamber foam builds up static, and a
dipole arm with no DC path to ground discharges straight into the analyzer
port. Cover both arms in heat-shrink, and give the feed a DC path from the
center pin to the shield with a shorted quarter-wave coax stub. Trim with both
in place.

\head{What Is Up to You}

This description is short on purpose: real projects arrive with a goal and a
deadline, not a procedure. What counts as ``tuned''? Which second frequency, and why that
one? Open stub or shorted, and how do you account for the velocity factor? Which
cuts, over what span, and how do you know your setup is not part of the
measurement? Make those calls, write down why you made them, and defend them in
the report.

A reference dipole with a published S$_{11}$ and pattern lives at the bench.
Measure it before you measure yours. If you cannot reproduce its numbers, the
problem is your setup, not your antenna.

\head{The Report}

One typed technical report, submitted as a PDF, with your measured data files
attached. There is no template and no page limit. Write it for an engineer who
was not in the lab with you: what you designed and why, how you built and
measured it, what you found, and what you make of it. Compare what you measured
against your 4NEC model. You are not required to make the numbers agree, but
you are required to know why they do not.

\head{Grading}

Show me each checkpoint at the start of the lesson listed.

\begin{center}
\begin{tabular}{|p{5.3in}|c|}
\hline
\textbf{What is assessed} & \textbf{Points} \\
\hline
\textbf{Checkpoints,} 5 points each. \textbf{L\cpOne:} the reference dipole
measured, and your trimmed dipole's S$_{11}$ and impedance at \fOne\ MHz.
\textbf{L\cpTwo:} co- and cross-polarized patterns at \fOne\ MHz, plus your
second frequency and the antenna's impedance there. \textbf{L\cpThree:} the
stub built, with its measured S$_{11}$. \textbf{L\cpFour:} co- and
cross-polarized patterns of the matched antenna. & 20 \\
\hline
\textbf{Analysis.} Measured against the 4NEC model at both frequencies, with every
disagreement explained or owned as unexplained. & 30 \\
\hline
\textbf{Measurements.} Impedance, S-parameters, and co- and cross-polarized
patterns, before and after the match, taken so that a reader can trust them. & 20 \\
\hline
\textbf{Design decisions.} The dipole, the trim, the second frequency, and the
stub, each with its reasoning stated. & 15 \\
\hline
\textbf{Report quality.} Organized, typed, every figure captioned with labeled
axes and units, data files included. & 15 \\
\hline
\end{tabular}
\end{center}

\textbf{Collaboration.} You may share chamber sessions and help each other take
data. The design, the data reduction, the analysis, and the report are
individual. Generative AI use is Level 3; document it as the syllabus requires.

\end{document}
